\begin{frame}{Industria Cloro-Álcali}
  \frametitle{Un Uso Residual en Proceso de Extinción}
  Históricamente, el mayor uso industrial del mercurio fue en la producción de cloro y sosa cáustica mediante celdas de mercurio.

  \begin{itemize}
    \item En este proceso, el mercurio actúa como cátodo líquido para producir una amalgama de sodio ($Na-Hg$) de alta pureza.
    \vspace{1cm}
    \item \textbf{Estado actual:} Esta tecnología se considera obsoleta. La mayoría de los países han prohibido su uso o han forzado la conversión de las plantas a tecnologías más seguras y eficientes, como las celdas de membrana.
    \vspace{1cm}
    \item Quedan muy pocas plantas operativas en el mundo y están bajo una estricta regulación para su cierre definitivo.
  \end{itemize}
\end{frame}

\begin{frame}{Iluminación Especializada}
  \frametitle{Aplicaciones de Nicho en Lámparas de Descarga}
  Cantidades muy pequeñas de mercurio son todavía esenciales para el funcionamiento de ciertos tipos de lámparas de alta eficiencia.

  \begin{itemize}
    \item El vapor de mercurio a baja presión es necesario para generar luz ultravioleta que luego excita el fósforo en \textbf{lámparas fluorescentes compactas (CFL)} y tubos lineales.
    \vspace{1cm}
    \item También se utiliza en lámparas de descarga de alta intensidad (HID) para alumbrado público y en \textbf{lámparas UV germicidas} para la esterilización de agua y aire.
    \vspace{1cm}
    \item \textbf{Innovación/Sustitución:} La tecnología \textbf{LED} está reemplazando rápidamente a las lámparas fluorescentes en casi todas las aplicaciones debido a su mayor eficiencia, vida útil y ausencia de mercurio.
  \end{itemize}
\end{frame}

\begin{frame}{Instrumentación Científica y de Control}
  \frametitle{Usos Altamente Especializados y de Referencia}
  Las propiedades físicas del mercurio lo mantienen como un material de referencia en aplicaciones científicas donde la precisión es crítica.

  \begin{itemize}
    \item \textbf{Porosimetría de mercurio:} Es una técnica estándar para medir la porosidad y la distribución del tamaño de los poros en materiales sólidos. No existe un sustituto directo que ofrezca la misma precisión.
    \vspace{1cm}
    \item \textbf{Electrodos de referencia:} Electrodos como el de calomelanos saturados ($Hg_2Cl_2$) son todavía patrones de referencia en electroquímica por su potencial estable y reproducible.
    \vspace{1cm}
    \item \textbf{Sustitución:} Se prioriza el uso de termómetros y manómetros digitales (basados en termistores y sensores piezorresistivos) en todas las aplicaciones no críticas.
  \end{itemize}
\end{frame}

\begin{frame}{El Futuro: Una Industria Libre de Mercurio}
  \frametitle{Principales Vías de Innovación Tecnológica}
  La investigación y el desarrollo se centran en eliminar los últimos bastiones del uso del mercurio.

  \begin{itemize}
    \item \textbf{Catálisis:} Desarrollo de catalizadores libres de mercurio para procesos como la producción de cloruro de vinilo (PVC), donde tradicionalmente se usaba cloruro de mercurio. Catalizadores a base de oro son una alternativa prometedora.
    \vspace{1cm}
    \item \textbf{Amalgamas dentales:} Aunque su uso ha decaído, la investigación se enfoca en mejorar la durabilidad y propiedades de las resinas compuestas y los ionómeros de vidrio para reemplazarlas por completo.
    \vspace{1cm}
    \item \textbf{Componentes electrónicos:} Desarrollo de interruptores y relés basados en tecnología de estado sólido (MEMS) para sustituir a los interruptores de mercurio en equipos especializados.
  \end{itemize}
\end{frame}

\begin{frame}{Puntos Importantes}
  \begin{itemize}
    \item El rol del mercurio en la industria moderna se ha reducido drásticamente y se limita a unas pocas \textbf{aplicaciones de nicho} donde aún no existen alternativas técnicamente viables.
    \vspace{1cm}
    \item El principal motor de la innovación no es encontrar nuevos usos para el mercurio, sino desarrollar \textbf{tecnologías de sustitución} eficientes, en línea con los objetivos del Convenio de Minamata.
    \vspace{1cm}
    \item El futuro industrial del mercurio no reside en su aplicación, sino en la gestión segura de sus residuos y en la \textbf{remediación} de la contaminación histórica.
  \end{itemize}
\end{frame}